\documentclass[11pt]{article}
\usepackage[margin=1in]{geometry}
\usepackage{amsmath,amssymb}
\usepackage{booktabs,longtable,array}
\usepackage{enumitem}
\usepackage{xcolor}
\usepackage{hyperref}
\usepackage{microtype}
\hypersetup{colorlinks=true,linkcolor=black,urlcolor=blue}
\newcommand{\codepath}[1]{\nolinkurl{#1}}
\title{\textbf{AutoIA: Roadmap from 0.21.0.dev1 to Version 1.0}\\
Release Blockers, Validation Gates, Performance Milestones, and Deferred Research}
\author{}
\date{August 2026}
\begin{document}
\maketitle

\begin{abstract}
AutoIA 0.21.0.dev1 already contains the central mathematical architecture needed
for an initial validated-execution release: machine-rigorous binary64 proof
kernels, polynomial and shared-generator representations, rigorously certified
smooth elementary functions and normal-distribution functions, semantic
parameter-dependent linear solves, path-aware scalar control flow, practical
subdivision, fast rigorous matrix reductions, and prepared repeated-query
execution with certificate restriction and multicore CPU batching.

The remaining question is no longer ``what else can AutoIA possibly support?''.
The correct Version-1 question is: ``what must be hardened, completed, measured,
and frozen so an outside engineer can install the package, understand its
supported semantics, obtain trustworthy results, and diagnose a failed proof?''
This roadmap therefore separates release blockers from useful follow-on research.
It intentionally avoids making GPU proof execution, a large semantic-primitive
zoo, or every NumPy operation prerequisites for Version 1.0.
\end{abstract}

\tableofcontents
\newpage

%===============================================================================
\section{Current baseline: what is already complete}
%===============================================================================

The current checkpoint is AutoIA \texttt{0.21.0.dev1}.  The final Step-9c tree
passes 223 regression/unit tests, builds an installable wheel, and passes a clean
wheel smoke test including generator-valued linear solve, exact/restricted
certificate cache reuse, and many-box bounds output.

The following major capabilities are already implemented:

\begin{enumerate}[leftmargin=*]
\item \textbf{Architecture and public failure taxonomy.}  Trusted numerical,
polynomial, enclosure, semantic, controller, and public API layers are separated;
proof/runtime policy state is context-local rather than process-global.
\item \textbf{Machine-rigorous binary64 reference backend.}  Environment audit,
outward scalar arithmetic, exact-rational regression, vectorized outward proof
reductions, and fail-closed numerical limits.
\item \textbf{Structured scalar representations.}  Interval, polynomial plus
remainder, shared-generator polynomial, and piecewise/path results.
\item \textbf{Smooth scalar functions.}  Validated \texttt{sqrt}, \texttt{exp},
\texttt{log}, \texttt{sin}, \texttt{cos}, constant real powers,
\texttt{erf}, normal PDF/CDF, plus common SciPy syntax compatibility.
\item \textbf{Ordinary NumPy engineering arrays.}  Object-array construction,
indexing/slicing, broadcasting, common reductions, dot/matmul, Euclidean/Frobenius
norms, and output shape retention for the documented subset.
\item \textbf{General semantic linear solve.}  Higher-order parameterized matrix
families, residual/contraction certification, preconditioner portfolio, multiple
right-hand sides, and generator-aware solve predictors.
\item \textbf{Uncertain scalar control flow.}  Path-aware \texttt{if}/\texttt{elif}
execution, exact simple boundaries, subdivision plus conservative path enumeration,
and fail-closed resource budgets.
\item \textbf{Practical adaptive subdivision.}  Cheap structural split score with
a realized-gain safeguard; expensive lookahead retained as a research oracle.
\item \textbf{Performance improvements.}  Rigorous vectorized/batched reductions,
proposal factorization reuse, sparse polynomial caches, and strong end-to-end
semantic-solve speedups without trusting BLAS as a proof oracle.
\item \textbf{Prepared repeated-query architecture.}  \texttt{prepare(model,...)},
exact cache, rigorous interior certificate restriction, serial/thread/process
many-box execution, and streaming bounds interfaces.
\end{enumerate}

This is enough functionality to justify moving toward an external release rather
than continuing an open-ended feature expansion.

%===============================================================================
\section{Definition of Version 1.0}
%===============================================================================

Version 1.0 should mean:

\begin{quote}
An outside user can install AutoIA on a supported platform, pass a documented
subset of ordinary engineering Python/NumPy code and interval input boxes, and
receive either (i) a machine-rigorous output certificate with documented
semantics or (ii) a typed fail-closed explanation.  The package has a frozen
public API, a reproducible independent validation suite, realistic performance
on representative engineering models, and enough documentation that a second
developer can safely extend the code without relying on undocumented project
history.
\end{quote}

Version 1.0 does \emph{not} need to mean:
\begin{itemize}[leftmargin=*]
\item every NumPy/SciPy operation is supported;
\item exact interval hulls are always obtained;
\item every failed proof can be rescued automatically;
\item GPU proof execution is available;
\item arbitrary ODE, nonlinear root, eigensolver, or optimization libraries are
semantic primitives;
\item all correlated physical input domains are represented intrinsically.
\end{itemize}

The release should be narrow enough to be defensible and broad enough to be
useful.

%===============================================================================
\section{Priority classes}
%===============================================================================

This roadmap uses four priority classes.

\begin{longtable}{p{0.16\textwidth}p{0.75\textwidth}}
\toprule
Priority & Meaning \\
\midrule
\endhead
P0 & Required before Version 1.0 because correctness, reproducibility, or basic external usability is not yet sufficiently mature without it. \\
P1 & Strongly desirable before Version 1.0, but may be deferred if it threatens release focus and the limitation can be clearly documented/fail-closed. \\
P2 & Natural post-1.0 capability/performance work enabled by the frozen architecture. \\
Research & Longer-term methodological extensions that should not delay the first stable release. \\
\bottomrule
\end{longtable}

%===============================================================================
\section{P0 blocker 1: resolve the centered-domain boundary issue cleanly}
%===============================================================================

The current polynomial input representation constructs a rigorous symmetric
centered domain that can extend slightly beyond the original user interval due to
outward rounding.  This is sound, but it can cause a valid physical domain such
as
\[
  x\in[0,4]
\]
to produce a represented polynomial argument range with a tiny negative lower
endpoint, causing \(\sqrt{x}\) or \(\log(x)\) to fail closed at an exact boundary.

This issue should be addressed before Version 1.0 because it is likely to confuse
outside users and because ad hoc clipping would be mathematically unacceptable.

\subsection{Recommended solution direction}

Introduce explicit physical-domain information distinct from the symmetric
polynomial coordinate box.  A validated scalar should be able to know both:
\begin{enumerate}[leftmargin=*]
\item the symmetric domain used to bound polynomial monomials; and
\item the exact/original physical interval information that may prove operation
specific domain constraints.
\end{enumerate}

For an operation such as \(\sqrt{x}\), AutoIA should prove nonnegativity of the
actual represented semantic quantity/domain rather than infer it by arbitrarily
clipping a polynomial range artifact.

\subsection{Acceptance tests}

P0 is complete when:
\begin{itemize}[leftmargin=*]
\item exact-boundary \(\sqrt{[0,b]}\) and appropriate \(\log\) boundary cases
have documented rigorous behavior;
\item the solution does not silently narrow a proof interval;
\item interior cases remain no-regression;
\item prepared certificate restriction respects the richer domain metadata;
\item boundary behavior is tested under both x86-64 and ARM64 supported
platforms.
\end{itemize}

%===============================================================================
\section{P0 blocker 2: batch-native polynomial/range performance (Step 9d)}
%===============================================================================

Step 9b moved the bottleneck away from rigorous matrix reductions.  Profiling now
identifies polynomial coefficient ranging and monotonicity-reduced Bernstein
transforms as dominant costs.  Step 9c also froze a many-box public API before
optimizing these internals.

The next performance implementation should therefore be \emph{batch-native}, not
an isolated one-box micro-optimization.

\subsection{Required work}

\begin{enumerate}[leftmargin=*]
\item Profile \texttt{range\_interval\_coefficient\_poly} and Bernstein conversion
on representative solve and non-solve models.
\item Replace repeated Python-object coefficient loops where possible by numeric
structure-of-arrays buffers.
\item Vectorize/batch monomial and coefficient transforms across outputs and, when
available, across many input boxes.
\item Preserve a slow scalar/reference path for regression and exact audit.
\item Revisit the Bernstein transform algorithmically; the earlier simple
precomputation attempt was measured and correctly reverted because it was slower.
\item Ensure any changed arithmetic order remains rigorously outward and covered
by independent exact/high-precision tests.
\end{enumerate}

\subsection{Release acceptance targets}

No single speed number should be a universal requirement, but Version 1 should
have:
\begin{itemize}[leftmargin=*]
\item an end-to-end benchmark suite with fixed representative models;
\item no major proof kernel dominated by avoidable Python loops at moderate array
sizes;
\item demonstrated useful scaling from one box to hundreds/thousands of boxes;
\item documented performance guidance for polynomial order, process workers,
cache use, and subdivision.
\end{itemize}

%===============================================================================
\section{P0 blocker 3: external-platform proof validation and CI}
%===============================================================================

The trusted kernel is audited at runtime, but Version 1 requires repeated release
validation on the platforms users are expected to run.

At minimum, establish continuous integration or reproducible release jobs for:
\begin{itemize}[leftmargin=*]
\item Linux x86-64;
\item macOS ARM64, important for current Apple Silicon engineering users;
\item preferably Windows x86-64 if Windows will be claimed as a supported
Version-1 platform.
\end{itemize}

Each supported platform should run:
\begin{enumerate}[leftmargin=*]
\item trusted environment audit;
\item exact-rational scalar arithmetic tests;
\item exact-rational reduction tests, including vectorized paths;
\item high-precision elementary/special-function tests;
\item near-singular semantic solve tests;
\item branch/path tests;
\item prepared certificate restriction tests;
\item wheel installation smoke tests.
\end{enumerate}

If a platform cannot satisfy the trusted binary64 assumptions, AutoIA should
explicitly refuse machine-rigorous execution there rather than weaken the audit.

%===============================================================================
\section{P0 blocker 4: external engineering validation corpus}
%===============================================================================

The current test suite is strong internally but heavily developed by the same
research effort that created the algorithms.  Version 1 should include models
written or selected independently of AutoIA implementation decisions.

\subsection{Recommended corpus}

Include at least:
\begin{itemize}[leftmargin=*]
\item mechanical spring/structural systems;
\item small thermal/fluid algebraic models;
\item parameterized linear systems with nonlinear coefficients;
\item models using normal CDF/PDF in reliability calculations;
\item branching engineering logic;
\item vector/matrix assembly and reductions;
\item repeated-query workloads with overlapping and unrelated interval boxes;
\item deliberately difficult near-singular and domain-boundary cases.
\end{itemize}

At least some models should come from people who did not participate in the core
AutoIA coding effort.

\subsection{What to measure}

For each model record:
\begin{itemize}[leftmargin=*]
\item whether unmodified user syntax is supported;
\item certificate success/failure and typed reason;
\item output widths and comparison to dense/high-precision sampling when feasible;
\item runtime and memory;
\item sensitivity to polynomial order and representation;
\item subdivision/path count;
\item repeated-query cache hit/restriction rate and multicore scaling.
\end{itemize}

Sampling is diagnostic, not the proof; its purpose is to identify unexpected
width inflation or implementation mistakes.

%===============================================================================
\section{P0 blocker 5: freeze the Version-1 supported-language contract}
%===============================================================================

The current subset is documented, but external testing will inevitably reveal
common idioms that are missing.  Before Version 1, make an explicit decision for
each common category: implement it or document a fail-closed exclusion.

Particular candidates to evaluate are:
\begin{itemize}[leftmargin=*]
\item scalar \texttt{abs}/\texttt{np.abs};
\item scalar and array \texttt{minimum}/\texttt{maximum} and \texttt{clip};
\item \texttt{np.where} with uncertain scalar/elementwise conditions;
\item common helper functions imported from neighboring Python modules;
\item local imports inside model functions;
\item common \texttt{einsum}/\texttt{tensordot} cases;
\item additional norms used in engineering code.
\end{itemize}

It is not necessary to implement all of these.  The release blocker is ambiguity.
A user should be able to determine from the manual whether a construct is
supported, and unsupported validated values must fail closed rather than be
silently coerced.

A short corpus study should drive this decision rather than adding operations
because they are easy in isolation.

%===============================================================================
\section{P0 blocker 6: API, packaging, and release engineering}
%===============================================================================

Before 1.0, freeze and test the public surface exposed from
\codepath{autoia/__init__.py}.  This includes:
\begin{itemize}[leftmargin=*]
\item \texttt{autoia};
\item \texttt{autoia\_subdivide};
\item \texttt{autoia\_adaptive\_subdivide};
\item \texttt{prepare} and \texttt{PreparedAutoIA};
\item interval/backend/policy/result classes;
\item special normal/error-function helpers;
\item public exception classes.
\end{itemize}

Required release-engineering tasks include:
\begin{enumerate}[leftmargin=*]
\item choose an actual software license; the current placeholder is not suitable
for outside release;
\item complete package metadata, project URLs, authorship/maintainer information,
and Python classifiers;
\item define supported Python/NumPy/SciPy version ranges and test them;
\item build wheels/sdist reproducibly;
\item add clean-environment installation tests;
\item ensure tests, historical oracle, and audit-only dependencies are not shipped
as runtime package code unless deliberately included;
\item add a concise public guarantee statement and supported-language reference;
\item add changelog/release notes and semantic-versioning rules.
\end{enumerate}

%===============================================================================
\section{P0 blocker 7: documentation and second-developer reproducibility}
%===============================================================================

The Version-1 documentation set should include, at minimum:
\begin{enumerate}[leftmargin=*]
\item a self-contained mathematical specification;
\item a theory-to-source implementation map;
\item a user/developer software manual;
\item the release roadmap and explicit deferred scope;
\item installation and several copyable engineering examples;
\item typed-error troubleshooting guidance;
\item instructions for running the independent audits and release tests.
\end{enumerate}

A practical acceptance criterion is that a technically capable second developer
who did not participate in the previous chat/research history can identify where
to add a new elementary function or semantic primitive and can explain which
calculations are proposals versus proof-critical calculations.

%===============================================================================
\section{P1: improve prepared many-box scaling beyond the current bounded scan}
%===============================================================================

Step 9c deliberately uses a simple bounded LRU cache and linear containment scan.
This is appropriate for freezing semantics, but a million-query workload may need
more sophisticated indexing.

Strong pre-1.0 improvements, if time permits, are:
\begin{itemize}[leftmargin=*]
\item spatial indexing of cached physical boxes;
\item persistent certificate stores for long-running jobs;
\item better policy for deciding when to certify a larger parent/master box versus
many individual children;
\item sharp restriction of pathwise/piecewise cached certificates;
\item metrics for cache coverage, reuse depth, and restriction cost.
\end{itemize}

These are P1 rather than P0 because the current behavior is already sound and
usable.  Version 1 can document that process-parallel fresh proofs are the fallback
for unrelated boxes.

%===============================================================================
\section{P1: numeric batch-buffer architecture}
%===============================================================================

The current tracing representation uses NumPy object arrays containing Python
validated scalar objects.  This is excellent for compatibility with ordinary
engineering code but not optimal for very large batched workloads or GPU
execution.

A desirable pre-1.0 internal refactor is a numeric structure-of-arrays buffer for
hot polynomial/range paths, with dimensions conceptually like
\[
  (N,\;n_{\mathrm{terms}},\;n_{\mathrm{outputs}}),
\]
where \(N\) is the number of simultaneously processed boxes.

This should remain an internal backend detail behind \texttt{PreparedAutoIA}; the
public API should not force users to know the buffer format.  If the work is too
large for the release schedule, the API can be frozen in 1.0 and the buffer
backend added afterward without breaking users.

%===============================================================================
\section{P1: generator growth and sound compression}
%===============================================================================

The shared-generator representation preserves useful provenance but can accumulate
many auxiliary IDs.  A naive merge can invent false correlation or destroy
important cancellation.

Any generator compression must therefore have an explicit inclusion contract.
Possible strategies include:
\begin{itemize}[leftmargin=*]
\item preserve a selected set of high-impact shared generators and convert the
rest to independent box error;
\item group generators only through a mathematically justified outer
approximation;
\item cap generator count with a documented monotone loss of correlation.
\end{itemize}

This is strongly desirable for large vector models, but not a correctness blocker
for Version 1 if the uncompressed mode remains usable and documented.

%===============================================================================
\section{P1: broader solve-family reuse}
%===============================================================================

Prepared repeated-query workloads may solve closely related matrix families many
times.  Additional safe proposal reuse could include:
\begin{itemize}[leftmargin=*]
\item previous-box preconditioners;
\item previous predictor coefficients;
\item reference matrix/factorization reuse when the exact point matrix is
unchanged;
\item shared matrix-family bounds across multiple right-hand sides;
\item warm starts across neighboring boxes.
\end{itemize}

Because these quantities are proposals, aggressive reuse is mathematically safe
provided each new box independently passes the residual/contraction certificate.
This is an attractive performance follow-on after Step 9d profiling.

%===============================================================================
\section{GPU roadmap: not a Version-1 proof blocker}
%===============================================================================

GPU capability should be pursued, but a full trusted GPU backend should not delay
Version 1.0.

\subsection{Stage G1: GPU proposal backend}

After numeric batch buffers exist, use a GPU for non-proof operations such as:
\begin{itemize}[leftmargin=*]
\item batched nominal/reference linear solves;
\item candidate preconditioners;
\item polynomial predictor coefficient generation;
\item controller estimates;
\item other large proposal-only matrix operations.
\end{itemize}
The CPU trusted backend then certifies the actual stored GPU proposal.  This gives
useful acceleration with minimal change to the trust boundary.

\subsection{Stage G2: experimental directed-rounding proof kernels}

A proof-producing GPU backend requires explicit control of directed rounding,
subnormal/flush-to-zero behavior, reduction ordering, and independent exact
validation.  Ordinary GPU NumPy analogues must not be assumed to reproduce the
CPU trusted semantics.

\subsection{Stage G3: optional trusted GPU release}

Only after cross-device audits and adversarial tests should GPU-produced bounds be
allowed to set \texttt{machine\_rigorous=True}.  Until then, GPU should be viewed
as a proposal accelerator.

%===============================================================================
\section{P2/post-1.0 semantic expansions}
%===============================================================================

Once the Version-1 architecture is stable, natural next semantic primitives
include:
\begin{itemize}[leftmargin=*]
\item nonlinear root systems using interval Newton/Krawczyk-style certificates;
\item eigensystem/singular-value bounds;
\item ODE/DAE integration with validated flow enclosures;
\item optimization semantic primitives;
\item additional probability distributions/special functions;
\item common sparse linear algebra;
\item richer non-box physical input domains.
\end{itemize}

These should be added only when there is an explicit mathematical semantics and
independent proof mechanism.  Their absence should not hold the first release
hostage.

%===============================================================================
\section{Suggested milestone sequence}
%===============================================================================

The exact version numbers are less important than the gates, but the following
sequence is recommended.

\begin{longtable}{p{0.18\textwidth}p{0.73\textwidth}}
\toprule
Milestone & Required outcome \\
\midrule
\endhead
M1: Step 9d & Batch-compatible polynomial coefficient/range/Bernstein optimization; profile after changes; preserve scalar oracle. \\
M2: Domain hardening & Resolve centered-domain boundary semantics without clipping; add boundary regression on supported platforms. \\
M3: External corpus & Run independent engineering models; decide the final common-language subset from actual usage; implement only high-value missing idioms. \\
M4: Platform/release hardening & CI on supported architectures, wheel/sdist pipeline, license/metadata, clean installation audits, version/dependency matrix. \\
M5: Release candidate & Freeze public API and guarantee statement; documentation review by a second developer/user; no new feature categories unless they fix a release blocker. \\
M6: Version 1.0 & Publish stable package with known limitations, benchmark/audit artifacts, and a post-1.0 research roadmap. \\
\bottomrule
\end{longtable}

%===============================================================================
\section{Quantitative release gates}
%===============================================================================

The project benefits from explicit go/no-go gates rather than a vague sense of
completeness.  Suggested gates are:

\begin{enumerate}[leftmargin=*]
\item \textbf{Regression gate:} all frozen tests pass on every supported release
platform.
\item \textbf{Trusted audit gate:} zero containment failures in exact-rational
scalar/reduction stress suites and independent high-precision elementary suites.
\item \textbf{External corpus gate:} all unsupported constructs fail with a typed
message; all returned certificates contain independent truth samples/reference
solutions; no silent numeric coercion is observed.
\item \textbf{Boundary gate:} documented \(\sqrt{\cdot}\)/\(\log\) physical-domain
boundary cases behave rigorously and predictably.
\item \textbf{Many-box gate:} prepared streaming can process a large synthetic
stream without unbounded memory growth; process mode scales materially on a
multicore machine; repeated/contained query reuse demonstrably avoids model
execution.
\item \textbf{Packaging gate:} clean wheel installation and smoke tests pass in
fresh environments for every supported Python version/platform.
\item \textbf{Documentation gate:} a second developer can build, test, modify,
and explain the proof boundary using only repository documentation.
\end{enumerate}

A million-box benchmark should be included before release, but the target should
be framed as throughput/memory/scaling characterization rather than an arbitrary
absolute runtime that depends strongly on model complexity and hardware.

%===============================================================================
\section{What should deliberately not be done before Version 1}
%===============================================================================

To prevent scope drift, the following should not become mandatory pre-1.0 work
unless external validation reveals a critical use case:

\begin{itemize}[leftmargin=*]
\item full rigorous GPU certification;
\item automatic recognition of arbitrary hand-coded linear/nonlinear algorithms;
\item a large zoo of alternative enclosure representations;
\item arbitrary correlated stochastic input objects;
\item full Python control-flow semantics;
\item general compiler/AST transformation of every helper module;
\item support for every SciPy distribution or special function;
\item global optimization or ODE certification merely because they are attractive
research topics.
\end{itemize}

The first release should establish credibility of the architecture and proof
contract.  New semantic domains can then build on that stable foundation.

%===============================================================================
\section{Immediate next action}
%===============================================================================

With Step 9c complete, the recommended immediate engineering action is:

\begin{quote}
Implement Step 9d by restructuring the dominant polynomial coefficient-ranging
and monotonicity-reduced Bernstein paths in a way that is compatible with the
prepared/many-box API and future numeric batch buffers.  Do not optimize the old
single-box object path in a way that would have to be discarded for batch/GPU
execution.
\end{quote}

In parallel, begin assembling the external engineering corpus and platform CI so
that correctness/usability findings can influence the remaining language and
domain work before the API is frozen.

%===============================================================================
\section{Final release perspective}
%===============================================================================

The core research question that motivated AutoIA has largely moved from
``can this architecture be made rigorous?'' to ``can it be made broad, fast, and
predictable enough for engineers outside the development effort?''  The major
mathematical pieces are now present.  The remaining Version-1 work should
therefore emphasize hardening, corpus-driven supported-language decisions,
performance where profiling shows real bottlenecks, repeated-query scalability,
and clear failure semantics.

A disciplined Version-1 scope is itself important research infrastructure.  It
creates a stable platform on which more ambitious semantic primitives, GPU proof
kernels, richer generator structures, and non-box uncertainty can be evaluated
without continually changing the foundational guarantee.

\end{document}
